\section{Conclusion}
\label{sec:conclusion}

We presented a large-scale study of the relationship between media bias and user interaction patterns on Men\'eame, a Spanish social news platform.
Our analysis of 14,995 articles with 13.2 million comments demonstrates that biased articles elicit measurably different engagement: higher karma entropy, more negative and angry comments, faster initial reactions, and steeper discourse degradation.
Community analysis reveals two structurally distinct media ecosystems among outlets, and user behavior analysis shows that moderate bias exposure correlates with the broadest information diets.

While interaction features alone achieve modest predictive performance (AUC $0.642$), they capture complementary signals to content-based approaches.
The consistency of these patterns across 16 years and multiple feature families suggests that crowd behavioral signals provide a robust, scalable layer for media bias monitoring.

Future work will explore: (1) combining interaction and content features in hybrid models, (2) obtaining comment threading information to analyze discussion depth and argumentation quality, (3) improving cross-lingual bias labeling to reduce dependence on small annotated corpora, and (4) extending the analysis to other social news platforms to test the generalizability of our findings.

\section*{Acknowledgments}

[CITE NEEDED: funding acknowledgments]

\section*{Ethical Considerations}

All user data was pseudonymized using SHA-256 hashing.
No attempt was made to deanonymize users.
The dataset contains only publicly available information from Men\'eame.
Bias labels are automatic and should not be interpreted as ground truth judgments about individual articles or outlets.
